\chapter{Math Prerequisites Refresher}
\label{ch:app_math}
\chaptermark{Math prerequisites}

This appendix is a refresher, not a textbook. If any of the concepts
below are unfamiliar, consult a dedicated source; the ``Further reading''
section at the end points to the standard references. The goal here is
to collect, in one place, the mathematical tools the body of the book
assumes --- so that when \cref{ch:kyle_gm} writes down a conditional
expectation, \cref{ch:mean_reversion_ou} invokes an Ornstein--Uhlenbeck
SDE, or \cref{ch:black_scholes} appeals to It\^o's lemma, the reader can
flip here and recover the definitions in a few seconds.

Nothing in this appendix is proved. Everything is stated precisely
enough to use. The rigour level matches that of a graduate-level
cheat sheet: state, don't derive.

\begin{backgroundbox}[How to use this appendix]
Read nothing here in sequence. Skip entire sections if you know the
material. Treat the appendix as a lookup table: when the body cites
``It\^o's lemma'' or ``Feynman--Kac'' or ``the spectral theorem for
symmetric matrices,'' come here, read the one paragraph that states the
result, and go back. If after reading the relevant paragraph the result
still does not click, the further-reading list at the end of this
appendix is your next stop --- not a re-read of this appendix.
\end{backgroundbox}

% =====================================================================
\section{Probability}
\label{sec:app_math_prob}

A probability model is a triple $(\Omega, \mathcal{F}, \Prob)$. The
\emph{sample space} $\Omega$ is the set of possible outcomes of a random
experiment. The \emph{$\sigma$-algebra} $\mathcal{F}$ is a collection
of subsets of $\Omega$ closed under complement and countable union; its
elements are \emph{events}. The \emph{probability measure}
$\Prob : \mathcal{F} \to [0,1]$ assigns a number to each event with
$\Prob(\Omega) = 1$ and $\Prob(\bigcup_i A_i) = \sum_i \Prob(A_i)$ for
disjoint $A_i$.

A \emph{random variable} is a measurable map $X : \Omega \to \R$.
Its \emph{distribution} is the pushforward measure
$\Prob \circ X^{-1}$ on $\R$; in practice one specifies it by a
\emph{cumulative distribution function} (CDF) $F_X(x) = \Prob(X \le x)$
or, when $X$ is absolutely continuous, by a \emph{probability density
function} (PDF) $f_X$ satisfying $F_X(x) = \int_{-\infty}^x f_X(u)\,du$.

\subsection{Expectation, variance, moments}

The \emph{expectation} (or mean) of $X$ is
\[
\E[X] = \int_\Omega X \, d\Prob
      = \int_\R x \, f_X(x)\, dx \quad \text{(continuous case)}.
\]
Expectation is linear: $\E[aX + bY] = a\E[X] + b\E[Y]$. The
\emph{variance} measures spread:
\[
\Var(X) = \E[(X - \E[X])^2] = \E[X^2] - (\E[X])^2,
\]
and the \emph{standard deviation} is $\sigma_X = \sqrt{\Var(X)}$.
Covariance and correlation of two variables are
\[
\Cov(X,Y) = \E[(X - \E X)(Y - \E Y)], \qquad
\Corr(X,Y) = \frac{\Cov(X,Y)}{\sigma_X \sigma_Y} \in [-1,1].
\]
The \emph{$k$-th moment} is $\E[X^k]$; the \emph{$k$-th central moment}
is $\E[(X - \E X)^k]$. Skewness and kurtosis are the standardised third
and fourth central moments:
\[
\operatorname{Skew}(X) = \frac{\E[(X-\mu)^3]}{\sigma^3}, \qquad
\operatorname{Kurt}(X) = \frac{\E[(X-\mu)^4]}{\sigma^4}.
\]
A Normal distribution has skewness $0$ and kurtosis $3$; \emph{excess
kurtosis} is $\operatorname{Kurt}(X) - 3$. Financial return
distributions (equities, crypto, most derivatives-linked series) have
excess kurtosis that is strongly positive --- ``fat tails.'' A naive
risk calculation that uses only mean and variance understates the
probability of large moves; see \cref{ch:kelly_risk}.

The \emph{moment generating function} (MGF) of $X$ is
$M_X(t) = \E[e^{tX}]$ when finite in a neighbourhood of $0$. MGFs, when
they exist, uniquely determine the distribution and linearise sums of
independent variables: $M_{X+Y}(t) = M_X(t)M_Y(t)$ for independent
$X, Y$. The \emph{characteristic function} $\varphi_X(t) = \E[e^{itX}]$
always exists and plays the same role --- it underpins the
Fourier-transform option-pricing methods used in \cref{ch:vol_surface}.

\begin{keyfactbox}[Linearity and variance identities]
\begin{itemize}
  \item $\E[aX + bY + c] = a\E[X] + b\E[Y] + c$ always.
  \item $\Var(aX + b) = a^2 \Var(X)$.
  \item $\Var(X + Y) = \Var(X) + \Var(Y) + 2\Cov(X,Y)$.
  \item If $X,Y$ are independent, $\Cov(X,Y) = 0$ and hence
        $\Var(X+Y) = \Var(X) + \Var(Y)$.
  \item Zero covariance does \emph{not} imply independence; it does for
        jointly Gaussian variables.
\end{itemize}
\end{keyfactbox}

\subsection{Conditional probability and Bayes' rule}

For events $A, B \in \mathcal{F}$ with $\Prob(B) > 0$,
\[
\Prob(A \mid B) = \frac{\Prob(A \cap B)}{\Prob(B)}.
\]
Events $A,B$ are \emph{independent} if $\Prob(A \cap B) = \Prob(A)\Prob(B)$,
equivalently $\Prob(A \mid B) = \Prob(A)$.

\emph{Bayes' rule} inverts the conditioning:
\[
\Prob(A \mid B) = \frac{\Prob(B \mid A)\,\Prob(A)}{\Prob(B)}
                = \frac{\Prob(B \mid A)\,\Prob(A)}
                       {\sum_i \Prob(B \mid A_i)\Prob(A_i)},
\]
where $\{A_i\}$ is a partition of $\Omega$. This is the engine behind
the Glosten--Milgrom market-maker update in \cref{ch:kyle_gm} and any
informed-trader posterior.

\emph{Example.} A coin is biased: $\Prob(\text{heads}) = 0.7$. We flip it
three times; what is the probability that the first flip was heads given
that exactly two of the three flips came up heads? Let $A = \{\text{first
is heads}\}$ and $B = \{\text{exactly two heads}\}$. Then
$\Prob(B) = \binom{3}{2}(0.7)^2(0.3) = 0.441$, while
$\Prob(A \cap B) = \Prob(\text{first H, one of next two H})
= 0.7 \cdot 2 \cdot 0.7 \cdot 0.3 = 0.294$. So
$\Prob(A \mid B) = 0.294 / 0.441 \approx 0.667$. Knowing that two of
three are heads has raised the posterior on the first flip from $0.7$
to $0.667$ --- wait, it \emph{lowered} it, because observing ``only
two heads'' is mild evidence against the first being heads. Compute
before intuiting.

\emph{Conditional expectation} of $X$ given a $\sigma$-algebra
$\mathcal{G} \subseteq \mathcal{F}$ is the $\mathcal{G}$-measurable
random variable $\E[X \mid \mathcal{G}]$ satisfying
$\E[\E[X \mid \mathcal{G}] \mathbf{1}_G] = \E[X \mathbf{1}_G]$ for all
$G \in \mathcal{G}$. Two facts used repeatedly in the body:
\begin{itemize}
  \item \emph{Tower property}: $\E[\E[X \mid \mathcal{G}]] = \E[X]$,
        and more generally $\E[\E[X \mid \mathcal{H}] \mid \mathcal{G}]
        = \E[X \mid \mathcal{G}]$ when $\mathcal{G} \subseteq \mathcal{H}$.
  \item \emph{Taking out what is known}: if $Y$ is $\mathcal{G}$-measurable,
        $\E[YX \mid \mathcal{G}] = Y\,\E[X \mid \mathcal{G}]$.
\end{itemize}
A \emph{martingale} (with respect to a filtration
$(\mathcal{F}_t)$) is a process $(M_t)$ with
(i) $M_t$ is $\mathcal{F}_t$-measurable, (ii) $\E[|M_t|] < \infty$, and
(iii) $\E[M_t \mid \mathcal{F}_s] = M_s$ for all $s \le t$.
A \emph{submartingale} replaces (iii) with $\ge$; a
\emph{supermartingale} with $\le$. Martingales are the natural model
for ``no drift conditional on information,'' hence their role in the
fundamental theorem of asset pricing (\cref{ch:black_scholes}).

Two martingale properties we rely on repeatedly:
\begin{itemize}
  \item \emph{Doob's optional stopping.} Under integrability
        conditions (for instance, if the stopping time $\tau$ is
        bounded), $\E[M_\tau] = \E[M_0]$. A strategy with bounded
        duration against a martingale price has expected payoff zero.
  \item \emph{Brownian motion and driftless It\^o integrals are
        martingales.} If $X_t$ is adapted and square-integrable, then
        $\int_0^t X_s\, dW_s$ is a martingale. This is why the
        risk-neutral expectation $\E_\Qm[\cdot]$ of a discounted payoff
        is itself the no-arbitrage price.
\end{itemize}

\subsection{Inequalities}

A handful of inequalities are used enough to deserve naming:
\begin{itemize}
  \item \emph{Markov:} for $X \ge 0$ and $a > 0$,
        $\Prob(X \ge a) \le \E[X]/a$.
  \item \emph{Chebyshev:} $\Prob(|X - \mu| \ge k\sigma) \le 1/k^2$.
        (Markov applied to $(X - \mu)^2$.)
  \item \emph{Cauchy--Schwarz:} $\E[XY]^2 \le \E[X^2]\E[Y^2]$. Equality
        iff $X = cY$ almost surely. This is what makes
        $|\Corr(X,Y)| \le 1$.
  \item \emph{Jensen:} for a convex function $\phi$ and integrable $X$,
        $\phi(\E[X]) \le \E[\phi(X)]$; for concave $\phi$ the inequality
        flips. Corollary: $\E[X^2] \ge (\E[X])^2$ (positive variance),
        and $\E[\log X] \le \log \E[X]$ (geometric mean $\le$
        arithmetic mean, the mathematical origin of the volatility
        drag $-\tfrac{1}{2}\sigma^2$ in log-returns and a running
        theme in \cref{ch:black_scholes} and \cref{ch:kelly_risk}).
\end{itemize}

\begin{intuitionbox}
Jensen is a statement about curvature. A convex-up function averages up
over any spread of its argument, so taking the expectation first loses
some upside relative to averaging the function values. For the
concave logarithm, this is why log-returns have a drag: you cannot
replace ``expected return'' with ``expected log return'' without losing
$\tfrac{1}{2}\sigma^2$, even in the absence of any trading cost.
\end{intuitionbox}

\subsection{Common distributions}

The distributions below are used throughout the book. The notation
$X \sim D$ means ``$X$ has distribution $D$.''

\begin{keyfactbox}[Distribution cheat sheet]
\begin{itemize}
  \item \textbf{Bernoulli}($p$): $\Prob(X=1)=p, \Prob(X=0)=1-p$.
        Mean $p$, variance $p(1-p)$.
  \item \textbf{Binomial}($n,p$): sum of $n$ i.i.d.\ Bernoulli($p$).
        Mean $np$, variance $np(1-p)$.
  \item \textbf{Poisson}($\lambda$): $\Prob(X=k)=e^{-\lambda}\lambda^k/k!$
        for $k=0,1,2,\dots$. Mean and variance both $\lambda$. Limit of
        Binomial($n, \lambda/n$) as $n\to\infty$; models independent
        arrivals in continuous time (order-book events, trade arrivals).
  \item \textbf{Exponential}($\lambda$): density $\lambda e^{-\lambda x}$
        for $x\ge 0$. Mean $1/\lambda$, variance $1/\lambda^2$. The
        inter-arrival time of a Poisson process. \emph{Memoryless}:
        $\Prob(X > s+t \mid X > s) = \Prob(X > t)$.
  \item \textbf{Normal} $\Normal(\mu, \sigma^2)$: density
        $\frac{1}{\sqrt{2\pi}\sigma} \exp\!\big(-(x-\mu)^2/(2\sigma^2)\big)$.
        Mean $\mu$, variance $\sigma^2$. Sums of independent Normals are
        Normal. Standard normal CDF denoted $\CDFN$, PDF $\pdfN$.
  \item \textbf{Log-Normal}: if $\log X \sim \Normal(\mu,\sigma^2)$ then
        $\E[X] = e^{\mu + \sigma^2/2}$ and
        $\Var(X) = (e^{\sigma^2}-1)\,e^{2\mu + \sigma^2}$.
        The Black--Scholes terminal stock price is log-normal
        (\cref{ch:black_scholes}).
\end{itemize}
\end{keyfactbox}

\subsection{Modes of convergence}

Four modes of convergence for a sequence of random variables $X_n$ to
a limit $X$ matter in what follows. From strongest to weakest:
\begin{itemize}
  \item \emph{Almost sure:} $X_n \xrightarrow{\text{a.s.}} X$ if
        $\Prob(\{\omega : X_n(\omega) \to X(\omega)\}) = 1$.
        ``All sample paths eventually settle.''
  \item \emph{In probability:} $X_n \xrightarrow{p} X$ if
        $\Prob(|X_n - X| > \varepsilon) \to 0$ for every
        $\varepsilon > 0$. Weaker than almost-sure convergence (which
        implies it).
  \item \emph{In $L^p$:} $X_n \xrightarrow{L^p} X$ if
        $\E[|X_n - X|^p] \to 0$. For $p = 2$ this is ``mean-square
        convergence''; implies convergence in probability.
  \item \emph{In distribution:} $X_n \xrightarrow{d} X$ if
        $F_{X_n}(x) \to F_X(x)$ at every continuity point of $F_X$.
        The weakest of the four.
\end{itemize}
The LLN below comes in two flavours: the \emph{weak law} (convergence
in probability) and the \emph{strong law} (almost-sure convergence).
The CLT is a statement about convergence in distribution.

\subsection{The multivariate normal}

A random vector $\bm{X} \in \R^n$ is \emph{multivariate normal} with
mean $\bm{\mu} \in \R^n$ and covariance $\Sigma \in \R^{n\times n}$
(PSD, often PD) if its density is
\[
f(\bm{x}) = \frac{1}{(2\pi)^{n/2}\, \det(\Sigma)^{1/2}}\,
  \exp\!\Big(-\tfrac{1}{2}(\bm{x} - \bm{\mu})^T \Sigma^{-1}(\bm{x}-\bm{\mu})\Big).
\]
We write $\bm{X} \sim \Normal(\bm{\mu}, \Sigma)$. Key properties:
\begin{itemize}
  \item \emph{Affine transforms stay normal:} if
        $\bm{X} \sim \Normal(\bm{\mu}, \Sigma)$ and $A$ is $m\times n$,
        $\bm{b} \in \R^m$, then
        $A\bm{X} + \bm{b} \sim \Normal(A\bm{\mu} + \bm{b}, A\Sigma A^T)$.
  \item \emph{Marginals are normal}, with the appropriate sub-vector of
        $\bm{\mu}$ and sub-matrix of $\Sigma$.
  \item \emph{Conditional distributions are normal.} Partitioning
        $\bm{X} = (\bm{X}_1, \bm{X}_2)^T$ with corresponding block
        structure $(\bm{\mu}_1, \bm{\mu}_2)$ and
        $\Sigma = \bigl(\begin{smallmatrix}\Sigma_{11} & \Sigma_{12}\\
                         \Sigma_{21} & \Sigma_{22}\end{smallmatrix}\bigr)$,
        \[
        \bm{X}_1 \mid \bm{X}_2 = \bm{x}_2
        \;\sim\; \Normal\!\Big(\bm{\mu}_1 + \Sigma_{12}\Sigma_{22}^{-1}(\bm{x}_2 - \bm{\mu}_2),\;
          \Sigma_{11} - \Sigma_{12}\Sigma_{22}^{-1}\Sigma_{21}\Big).
        \]
        This formula is the entirety of the Glosten--Milgrom
        conditional-update in \cref{ch:kyle_gm}, Kalman filtering
        for latent-state estimation, and factor-model residual
        decomposition.
  \item \emph{Sampling.} To sample $\bm{X} \sim \Normal(\bm{0}, \Sigma)$,
        Cholesky-factor $\Sigma = LL^T$ and set $\bm{X} = L\bm{Z}$ for
        $\bm{Z} \sim \Normal(\bm{0}, I)$. The standard recipe for
        simulating correlated risk factors.
\end{itemize}

\subsection{Limit theorems}

Two statements dominate classical probability, and both underpin
backtesting and Monte Carlo in this book.

\emph{Law of Large Numbers.} Let $X_1, X_2, \dots$ be i.i.d.\ with finite
mean $\mu = \E[X_1]$. Then the sample average converges to the mean:
\[
\bar X_n \;=\; \frac{1}{n}\sum_{i=1}^n X_i \;\xrightarrow{\text{a.s.}}\; \mu
\quad \text{as } n \to \infty.
\]

\emph{Central Limit Theorem.} With the same assumptions and additionally
$\Var(X_1) = \sigma^2 < \infty$,
\[
\sqrt{n}\,(\bar X_n - \mu) \;\xrightarrow{d}\; \Normal(0, \sigma^2)
\quad \text{as } n \to \infty.
\]
In words: sample averages are approximately Normal, with standard error
$\sigma / \sqrt{n}$.

\begin{keyfactbox}[Rate of convergence in CLT]
The CLT gives the asymptotic \emph{scaling} --- $1/\sqrt{n}$ --- but not
a distribution at finite $n$. The Berry--Esseen theorem sharpens this:
the Kolmogorov distance between the CDF of $\sqrt{n}(\bar X_n - \mu)/\sigma$
and the standard normal is bounded by $C \E[|X_1 - \mu|^3]/(\sigma^3 \sqrt{n})$,
with $C$ a universal constant (${\le 0.4748}$ in the i.i.d.\ case).
For backtest Sharpe-ratio inference with daily returns and $n \gtrsim 500$
observations, the Gaussian approximation is usually adequate; for
intraday strategies with hundreds of thousands of trades it is almost
exact.
\end{keyfactbox}

\begin{pitfallbox}[Finite-variance assumption matters]
The CLT requires $\sigma^2 < \infty$. Financial returns often exhibit
heavy tails (kurtosis $\gg 3$); in the extreme (Pareto tail index
$\alpha < 2$) the variance is infinite and the CLT fails --- sample
averages converge to a non-Gaussian stable law, or do not converge at
all. Backtests that assume Gaussian standard errors on Sharpe ratios
(\cref{ch:kelly_risk}) silently mis-size their confidence intervals
when returns are that heavy-tailed.
\end{pitfallbox}

% =====================================================================
\section{Linear algebra}
\label{sec:app_math_linalg}

Vectors $\bm{x} \in \R^n$ are columns; matrices $A \in \R^{m\times n}$
act on them by left multiplication. The \emph{inner product} is
$\langle \bm{x}, \bm{y}\rangle = \bm{x}^T \bm{y} = \sum_i x_i y_i$ and
the Euclidean norm is $\|\bm{x}\| = \sqrt{\langle \bm{x},\bm{x}\rangle}$.

\subsection{Matrix arithmetic}

Matrix multiplication is associative but not commutative: in general
$AB \ne BA$. The \emph{transpose} $A^T$ flips rows and columns;
$(AB)^T = B^T A^T$ and $(AB)^{-1} = B^{-1}A^{-1}$ when invertible.
The \emph{trace} $\operatorname{tr}(A) = \sum_i A_{ii}$ is linear and
cyclic: $\operatorname{tr}(AB) = \operatorname{tr}(BA)$,
$\operatorname{tr}(ABC) = \operatorname{tr}(BCA) = \operatorname{tr}(CAB)$.

Three useful factorisations beyond eigen/SVD:
\begin{itemize}
  \item \textbf{LU}: $PA = LU$ where $P$ is a permutation, $L$ lower
        triangular (unit diagonal), $U$ upper triangular. Standard
        general-purpose linear-system solver. $\mathcal{O}(n^3)$ to
        compute, $\mathcal{O}(n^2)$ to back-substitute.
  \item \textbf{Cholesky}: $A = LL^T$ when $A$ is symmetric PD. About
        twice as fast as LU and numerically stable. The default for
        covariance matrices and normal equations.
  \item \textbf{QR}: $A = QR$ where $Q$ has orthonormal columns and $R$
        is upper triangular. The preferred way to solve
        least-squares --- avoids forming the ill-conditioned $A^T A$.
\end{itemize}

\subsection{Rank, inverse, determinant}

The \emph{rank} of $A$ is the dimension of its column space (equivalently
its row space). $A \in \R^{n\times n}$ is \emph{invertible} iff
$\operatorname{rank}(A) = n$, iff $\det(A) \ne 0$, iff
$A\bm{x} = \bm{0}$ has only the trivial solution.

The \emph{determinant} satisfies $\det(AB) = \det(A)\det(B)$ and
$\det(A^{-1}) = 1/\det(A)$. Geometrically, $|\det(A)|$ is the factor by
which $A$ scales $n$-dimensional volume.

The \emph{inverse}, when it exists, satisfies $AA^{-1}=A^{-1}A = I$. For
$2\times 2$ matrices,
\[
\begin{pmatrix} a & b \\ c & d \end{pmatrix}^{-1}
= \frac{1}{ad-bc}\begin{pmatrix} d & -b \\ -c & a \end{pmatrix}.
\]
Avoid inverting explicitly in code --- solve the linear system
$A\bm{x} = \bm{b}$ directly (LU, QR, or tridiagonal; see
\cref{ch:numerical_methods}).

\subsection{Eigenvalues and eigenvectors}

A scalar $\lambda \in \mathbb{C}$ and non-zero $\bm{v} \in \mathbb{C}^n$
are an \emph{eigenvalue}--\emph{eigenvector} pair of $A \in \R^{n\times n}$
if $A \bm{v} = \lambda \bm{v}$. Equivalently, $\lambda$ is a root of the
\emph{characteristic polynomial} $\det(A - \lambda I) = 0$. An
$n\times n$ matrix has $n$ eigenvalues counted with multiplicity.

\emph{Trace and determinant relations}: $\operatorname{tr}(A) = \sum_i \lambda_i$
and $\det(A) = \prod_i \lambda_i$.

\begin{keyfactbox}[Spectral theorem for symmetric matrices]
If $A \in \R^{n\times n}$ is symmetric ($A = A^T$), then:
\begin{enumerate}
  \item all eigenvalues $\lambda_i$ are real;
  \item eigenvectors corresponding to distinct eigenvalues are
        orthogonal, and one can choose an orthonormal basis of
        eigenvectors;
  \item consequently $A = Q \Lambda Q^T$ where $Q$ is orthogonal
        ($Q^T Q = I$) and $\Lambda = \operatorname{diag}(\lambda_1,\dots,\lambda_n)$.
\end{enumerate}
This is the \emph{eigendecomposition}. It is the mathematical skeleton
behind PCA, portfolio variance minimisation, and the asymptotic
distribution of OU processes in \cref{ch:mean_reversion_ou}.
\end{keyfactbox}

\emph{Diagonalisation.} If $A$ has $n$ linearly independent
eigenvectors (always the case for symmetric $A$, sometimes for
non-symmetric), then $A = P D P^{-1}$ with $D$ diagonal. Powers are
trivial: $A^k = P D^k P^{-1}$. Matrix exponentials are
$e^A = P e^D P^{-1}$, and $e^D$ is just the diagonal of scalar
exponentials. This is what makes the solution of a linear
homogeneous ODE $\dot {\bm x} = A \bm x$ trivial in the
diagonalisable case: $\bm x(t) = e^{At} \bm x(0) = P e^{Dt} P^{-1}
\bm x(0)$.

\emph{Worked example.} Consider
\[
A = \begin{pmatrix} 2 & 1 \\ 1 & 2 \end{pmatrix}.
\]
The characteristic polynomial is
$\det(A - \lambda I) = (2-\lambda)^2 - 1 = \lambda^2 - 4\lambda + 3$,
with roots $\lambda_1 = 3$ and $\lambda_2 = 1$. Eigenvectors:
for $\lambda = 3$, $(A - 3I)\bm{v} = 0$ gives $\bm{v} \propto (1,1)^T$;
for $\lambda = 1$, $\bm{v} \propto (1,-1)^T$. After normalisation the
orthogonal matrix of eigenvectors is $Q = \tfrac{1}{\sqrt{2}}\begin{pmatrix}
1 & 1 \\ 1 & -1 \end{pmatrix}$ and $A = Q \operatorname{diag}(3,1) Q^T$.
Sanity check: $\operatorname{tr}(A) = 4 = 3 + 1$ and $\det(A) = 3 =
3 \cdot 1$.

A symmetric matrix $A$ is \emph{positive semi-definite} (PSD) if
$\bm{x}^T A \bm{x} \ge 0$ for all $\bm{x}$, equivalently all
eigenvalues $\ge 0$. It is \emph{positive definite} (PD) if the
inequality is strict for $\bm{x} \ne \bm{0}$. Covariance matrices are
always PSD. A PD matrix has a \emph{Cholesky factorisation}
$A = LL^T$ with $L$ lower triangular --- the default way to sample
from a multivariate Normal $\Normal(\bm{0}, A)$ in code.

\subsection{Singular value decomposition}

For any $A \in \R^{m\times n}$ (not necessarily square), the
\emph{singular value decomposition} (SVD) is
\[
A = U \Sigma V^T
\]
where $U \in \R^{m\times m}$ and $V \in \R^{n\times n}$ are orthogonal,
and $\Sigma \in \R^{m\times n}$ is diagonal with non-negative entries
$\sigma_1 \ge \sigma_2 \ge \cdots \ge \sigma_r > 0$ (the \emph{singular
values}), where $r = \operatorname{rank}(A)$. The columns of $U$ and $V$
are the \emph{left} and \emph{right singular vectors}.

The SVD generalises eigendecomposition to rectangular matrices and to
non-symmetric square matrices. Numerically it is the most stable way to
diagnose near-singularity (look at the smallest singular value relative
to the largest; the ratio $\sigma_1/\sigma_r$ is the \emph{condition
number}). A condition number above, say, $10^8$ in double precision is
a strong signal that small perturbations of the inputs produce large
perturbations of the outputs --- a common failure mode when regressing
almost-collinear features (two highly-correlated price series, the
classic pairs-trading hazard of \cref{ch:pairs_cointegration}).

The SVD also gives the Moore--Penrose \emph{pseudoinverse}
$A^+ = V \Sigma^+ U^T$, where $\Sigma^+$ is formed by inverting the
non-zero singular values and transposing. The least-squares solution
of $A\bm{x} = \bm{b}$ (the projection of $\bm{b}$ onto the column space
of $A$, minimising $\|A\bm{x} - \bm{b}\|$) is $\bm{x}^* = A^+ \bm{b}$.
When $A$ has full column rank, $A^+ = (A^T A)^{-1}A^T$ --- the formula
behind every ordinary-least-squares regression, including the hedge-ratio
estimation in \cref{ch:pairs_cointegration}.

\subsection{PCA as eigendecomposition of the covariance}

Given data matrix $X \in \R^{n \times p}$ (rows are observations,
columns are features) with column means subtracted, the sample
covariance is $\hat\Sigma = \frac{1}{n-1} X^T X \in \R^{p\times p}$.
$\hat\Sigma$ is symmetric PSD, so by the spectral theorem
$\hat\Sigma = Q \Lambda Q^T$.

\emph{Principal component analysis} (PCA) is: order the eigenvalues
$\lambda_1 \ge \lambda_2 \ge \cdots \ge \lambda_p$; the $k$-th
\emph{principal component} is the direction $\bm{q}_k$ (the $k$-th column
of $Q$); the variance explained by the $k$-th component is $\lambda_k$
and the fraction of total variance is $\lambda_k / \sum_j \lambda_j$.
Keeping the top $k < p$ components gives the best rank-$k$ approximation
to $X$ in Frobenius norm (Eckart--Young).

Equivalently, the PCA directions are the right singular vectors of the
(centred) data matrix $X$, and the singular values satisfy
$\sigma_i^2 = (n-1)\lambda_i$.

\begin{intuitionbox}
PCA rotates the feature axes to line up with the directions of greatest
variance. In trading: run PCA on a panel of stock returns; the first
component is typically ``the market,'' the second is ``long one sector,
short another,'' and so on. The eigenvectors give the portfolio weights,
the eigenvalues give the variance of each resulting portfolio. This is
the foundation of statistical-arbitrage factor models in
\cref{ch:pairs_cointegration}.
\end{intuitionbox}

\begin{pitfallbox}[Scale matters for PCA]
PCA on raw data is dominated by whichever feature has the largest
variance, regardless of signal. Standardise (subtract mean, divide by
standard deviation) each column before running PCA unless the features
are already on a comparable scale. For returns panels this is usually
automatic; for mixed-unit features (a price and a volume, say) it is
mandatory.
\end{pitfallbox}

% =====================================================================
\section{Stochastic calculus sketch}
\label{sec:app_math_stoch}

This section is a one-page tour of the machinery used in
\cref{ch:kyle_gm}, \cref{ch:mean_reversion_ou}, \cref{ch:black_scholes},
\cref{ch:almgren_chriss}, and \cref{ch:numerical_methods}. For actual
study, use \citet{oksendal2003stochastic} or \citet{shreve2004stochastic}.

\subsection{Brownian motion}

A \emph{standard Brownian motion} (or \emph{Wiener process}) $W_t$ is a
continuous-time stochastic process with:
\begin{enumerate}
  \item $W_0 = 0$;
  \item \emph{continuous paths} $t \mapsto W_t$ (almost surely);
  \item \emph{independent increments}: for $s < t$, $W_t - W_s$ is
        independent of $(W_u : u \le s)$;
  \item \emph{Gaussian increments}: $W_t - W_s \sim \Normal(0, t-s)$.
\end{enumerate}
Brownian motion is the unique (up to scaling) process satisfying
these properties. Its paths are almost surely nowhere differentiable,
but they have a very precise second-order behaviour, quantified by
quadratic variation.

\subsection{Quadratic variation and the rule $(dW)^2 = dt$}

The \emph{quadratic variation} of $W$ on $[0,t]$ is
\[
[W]_t \;=\; \lim_{\|\pi\|\to 0}\; \sum_{k} (W_{t_{k+1}} - W_{t_k})^2
\;=\; t
\]
almost surely, where the limit is over partitions
$\pi = \{0 = t_0 < t_1 < \cdots < t_n = t\}$ with mesh $\|\pi\| \to 0$.
This is the content of the heuristic
\[
(dW_t)^2 = dt, \qquad dW_t \, dt = 0, \qquad (dt)^2 = 0,
\]
which is the arithmetic rule that makes all of stochastic calculus
``work.'' Squared Brownian increments do not vanish in the limit the way
$(dt)^2$ does --- they accumulate deterministically to $dt$.

\subsection{It\^o's lemma}

For a smooth function $f(W_t, t)$, the ordinary chain rule would give
$df = \partial_t f \, dt + \partial_W f \, dW_t$. The correction is the
$(dW)^2 = dt$ term of quadratic variation:

\begin{keyfactbox}[It\^o's lemma (scalar form)]
For $f \in C^{1,2}(\R\times [0,T])$ and $W_t$ a standard Brownian motion,
\[
df(W_t, t) \;=\; \partial_t f(W_t, t)\, dt
              \;+\; \partial_W f(W_t, t)\, dW_t
              \;+\; \tfrac{1}{2}\, \partial_{WW} f(W_t, t)\, dt.
\]
More generally, if $X_t$ satisfies $dX_t = \mu_t\, dt + \sigma_t\, dW_t$,
then for $f(X_t, t)$,
\[
df = \Big(\partial_t f + \mu_t\, \partial_x f
         + \tfrac{1}{2}\sigma_t^2\, \partial_{xx} f\Big)\, dt
     \;+\; \sigma_t\, \partial_x f\, dW_t.
\]
\end{keyfactbox}

The extra $\tfrac{1}{2}\partial_{xx} f$ term is the \emph{It\^o
correction}: it comes from the fact that $W$ has non-trivial
second-order variation. For deterministic $x$, it vanishes
($\sigma_t \equiv 0$) and one recovers ordinary calculus.

\subsection{Stochastic differential equations}

A \emph{stochastic differential equation} (SDE) is an equation of the
form
\[
dX_t = \mu(X_t, t)\, dt + \sigma(X_t, t)\, dW_t, \qquad X_0 = x_0,
\]
interpreted as the integral equation
\[
X_t = x_0 + \int_0^t \mu(X_s, s)\, ds + \int_0^t \sigma(X_s, s)\, dW_s,
\]
where the second integral is an \emph{It\^o stochastic integral}. Under
mild Lipschitz and linear-growth conditions on $\mu$ and $\sigma$, the
SDE admits a unique strong solution.

An SDE is a probabilistic generalisation of an ODE in which the
velocity field acquires a random component. If $\sigma \equiv 0$, the
SDE reduces to the ODE $\dot X_t = \mu(X_t, t)$. The new ingredient is
that the Brownian differential $dW_t$ is $\mathcal{O}(\sqrt{dt})$, not
$\mathcal{O}(dt)$ --- which is why the second-order term in It\^o's
lemma is not negligible.

\emph{Multi-dimensional It\^o.} For $X_t \in \R^n$ with
$dX_t = \bm{\mu}(X_t, t)\, dt + \Sigma(X_t, t)\, dW_t$ (here
$W_t \in \R^d$ and $\Sigma \in \R^{n\times d}$), and $f(\bm{x}, t) \in \R$,
\[
df = \Big(\partial_t f + \bm{\mu}^T \nabla f
   + \tfrac{1}{2}\operatorname{tr}(\Sigma \Sigma^T \nabla^2 f)\Big) dt
   + (\nabla f)^T \Sigma\, dW_t.
\]
The generator becomes
$L f = \bm{\mu}^T \nabla f + \tfrac{1}{2}\operatorname{tr}(\Sigma\Sigma^T \nabla^2 f)$,
a linear second-order elliptic operator; this is what appears in
multi-asset Black--Scholes and in basket/spread options pricing
(\cref{ch:basket_etf_arb}).

Three SDEs recur throughout the body:
\begin{itemize}
  \item \emph{Arithmetic Brownian motion} $dX_t = \mu\, dt + \sigma\, dW_t$.
        Solution $X_t = x_0 + \mu t + \sigma W_t$. Normal at every
        horizon. Used as a toy model and for latent fair values in
        Kyle-style microstructure (\cref{ch:kyle_gm}).
  \item \emph{Geometric Brownian motion} $dS_t = \mu S_t\, dt + \sigma S_t\, dW_t$.
        By It\^o applied to $f = \log S_t$, the solution is
        $S_t = S_0 \exp\!\big((\mu - \tfrac{1}{2}\sigma^2)t + \sigma W_t\big)$,
        which is log-normal. Black--Scholes assumes GBM
        (\cref{ch:black_scholes}).
  \item \emph{Ornstein--Uhlenbeck process} $dX_t = \theta(\mu - X_t)\,dt + \sigma\,dW_t$.
        Mean-reverts to $\mu$ with speed $\theta$; stationary
        distribution $\Normal(\mu, \sigma^2/(2\theta))$; half-life
        $\log 2 / \theta$. The workhorse model of
        \cref{ch:mean_reversion_ou} and
        \cref{ch:pairs_cointegration}.
\end{itemize}

\emph{Solving the OU SDE explicitly.} Apply It\^o to
$f(X_t, t) = X_t e^{\theta t}$. The partial derivatives are
$\partial_t f = \theta X_t e^{\theta t}$ and
$\partial_x f = e^{\theta t}$, with $\partial_{xx}f = 0$. Hence
\[
df = \theta X_t e^{\theta t} dt + e^{\theta t} dX_t
   = \theta X_t e^{\theta t} dt + e^{\theta t}\big[\theta(\mu - X_t)dt + \sigma dW_t\big]
   = \theta\mu e^{\theta t}\, dt + \sigma e^{\theta t}\, dW_t.
\]
Integrating from $0$ to $t$,
\[
X_t e^{\theta t} - X_0 = \mu(e^{\theta t} - 1) + \sigma \int_0^t e^{\theta s}\, dW_s,
\]
so
\[
X_t = X_0 e^{-\theta t} + \mu(1 - e^{-\theta t}) + \sigma \int_0^t e^{-\theta(t-s)} dW_s.
\]
The last term is a Gaussian with mean zero and variance
$\sigma^2 \int_0^t e^{-2\theta(t-s)} ds = \sigma^2 (1-e^{-2\theta t})/(2\theta)$,
which tends to $\sigma^2/(2\theta)$ as $t \to \infty$ --- recovering the
stationary distribution. This is the exact solution used to simulate OU
processes without discretisation error in
\cref{ch:mean_reversion_ou}.

\subsection{Feynman--Kac formula}

The Feynman--Kac formula is the bridge between PDEs and expectations
under a diffusion. It is what lets Black--Scholes be stated either as a
PDE (\cref{ch:black_scholes}, equation for option price) or as a
discounted expectation under the risk-neutral measure --- both views are
the same equation read in two directions.

\begin{keyfactbox}[Feynman--Kac (statement only)]
Let $X_t$ satisfy $dX_t = \mu(X_t, t)\, dt + \sigma(X_t, t)\, dW_t$, let
$L$ be the infinitesimal generator
\[
L = \mu(x, t)\,\partial_x + \tfrac{1}{2}\sigma(x,t)^2\,\partial_{xx},
\]
and let $r(x,t), g(x)$ be (bounded, sufficiently smooth) functions. If
$u(x,t)$ solves the backward PDE
\[
\partial_t u(x,t) + Lu(x,t) - r(x,t)\,u(x,t) = 0,
\qquad u(x,T) = g(x),
\]
then, subject to standard integrability conditions,
\[
u(x,t) \;=\; \E\!\left[\exp\!\Big(-\!\int_t^T r(X_s, s)\,ds\Big)\,
                         g(X_T) \;\Big|\; X_t = x\right].
\]
\end{keyfactbox}

The takeaway: \emph{backward PDEs with terminal conditions are
expectations of terminal payoffs under the diffusion, optionally
discounted}. The Black--Scholes PDE with payoff $g(x) = (x - \Kstrike)^+$
and discount $r$ is exactly such an expectation under the risk-neutral
dynamics --- giving the closed-form call price of \cref{ch:black_scholes}.

\subsection{Girsanov and change of measure}

Pricing under the risk-neutral measure $\Qm$ (\cref{ch:black_scholes})
requires a tool to move between the physical measure $\Prob$ (under
which the stock drifts at some real-world rate $\mu$) and $\Qm$ (under
which it drifts at the risk-free rate $r$).

\emph{Girsanov's theorem}, in its most-used form, says: let $W_t$ be a
$\Prob$-Brownian motion and $\theta_t$ a suitably integrable adapted
process. Define
\[
Z_t = \exp\!\Big(-\!\int_0^t \theta_s\, dW_s - \tfrac{1}{2}\int_0^t \theta_s^2\, ds\Big),
\]
and a new measure $\Qm$ by $d\Qm/d\Prob = Z_T$. Then
$\widetilde W_t = W_t + \int_0^t \theta_s\, ds$ is a $\Qm$-Brownian
motion. In other words: changing the drift of a diffusion by $\theta_t$
is equivalent to reweighting the probability measure by $Z_t$.

The practical content: under the risk-neutral measure $\Qm$, all
discounted tradeable asset prices are martingales. The
\emph{fundamental theorem of asset pricing} states that absence of
arbitrage is equivalent to the existence of such a measure; completeness
of the market is equivalent to its uniqueness. Everything that
follows in Black--Scholes, Heston, local vol, and the LIBOR Market
Model is Girsanov applied cleverly.

\begin{pitfallbox}[$\Prob$ versus $\Qm$]
Parameters calibrated from historical data (realised vol, realised
drift) live under $\Prob$; parameters backed out from option prices
(implied vol, implied forward) live under $\Qm$. You cannot mix them
in a single formula without explicitly writing down the Radon--Nikodym
derivative. A common error in junior-strat code is to price an option
with a risk-neutral BS formula while plugging in the historical drift
of the underlying --- the wrong measure, and the result is neither a
fair value nor a forecast.
\end{pitfallbox}

% =====================================================================
\section{Stochastic control sketch}
\label{sec:app_math_control}

Stochastic control asks: given a controlled diffusion and a cost or
reward over $[0,T]$, what control $\alpha_t$ (chosen dynamically, using
only information available at time $t$) maximises expected reward?
This is the framework of \cref{ch:market_making} (Avellaneda--Stoikov
market-making) and \cref{ch:almgren_chriss} (optimal execution).

\subsection{Setup and the value function}

Given a controlled SDE
\[
dX_t = \mu(X_t, \alpha_t)\,dt + \sigma(X_t, \alpha_t)\,dW_t,
\]
a running reward $r(x, \alpha)$, and a terminal reward $g(x)$, define
the \emph{value function}
\[
V(x, t) \;=\; \sup_{\alpha} \;
  \E\!\left[\int_t^T r(X_s, \alpha_s)\,ds + g(X_T) \;\Big|\; X_t = x\right],
\]
where the supremum is over \emph{admissible} (progressively measurable,
integrable) controls. $V(x, t)$ is ``the best expected future reward
achievable starting from state $x$ at time $t$ and playing optimally
from there.''

The guiding principle is \emph{dynamic programming}: an optimal
trajectory is optimal over any sub-interval. Formally (Bellman principle),
for any stopping time $\tau \in [t, T]$,
\[
V(x,t) = \sup_\alpha \E\!\left[
  \int_t^\tau r\,ds + V(X_\tau, \tau) \,\Big|\, X_t = x
\right].
\]

\subsection{Hamilton--Jacobi--Bellman equation}

Passing Bellman's principle to the infinitesimal limit $\tau \downarrow t$
and applying It\^o gives the \emph{Hamilton--Jacobi--Bellman} (HJB)
partial differential equation:

\begin{keyfactbox}[HJB equation (statement only)]
Under standard regularity assumptions, the value function $V(x,t)$
satisfies
\[
\partial_t V(x,t) + \sup_{\alpha}
  \Big\{ L^\alpha V(x,t) + r(x,\alpha) \Big\} = 0,
\qquad V(x,T) = g(x),
\]
where $L^\alpha$ is the generator of the controlled diffusion under
control $\alpha$:
\[
L^\alpha f = \mu(x,\alpha)\,\partial_x f + \tfrac{1}{2}\sigma(x,\alpha)^2\,\partial_{xx}f.
\]
The optimal feedback control at $(x,t)$ is any $\alpha^*(x,t)$ achieving
the supremum inside the braces.
\end{keyfactbox}

HJB is the stochastic-control analogue of Bellman's equation in
discrete-time dynamic programming. Solving it --- either in closed form
(Avellaneda--Stoikov, Almgren--Chriss, Merton's portfolio problem), by
reduction to a linear PDE via a clever transformation, or numerically on
a grid --- is how we pin down optimal quoting (\cref{ch:market_making})
and optimal execution (\cref{ch:almgren_chriss}).

A first-pass reader should leave this section with three facts: (i) the
value function captures best-attainable future reward; (ii) HJB is a
backward PDE with a sup over the control inside; (iii) the optimal
control is whichever argument attains that sup. The body of the book
applies HJB to two cases where it can be solved in closed form and one
where it is solved numerically on a finite grid.

\subsection{A tiny worked HJB: Merton's portfolio problem}

To make the machinery concrete, consider Merton's canonical problem:
invest wealth $X_t$ in a risky asset (GBM with drift $\mu$, vol
$\sigma$) and a risk-free asset (rate $r$), with control $\alpha_t$ =
fraction of wealth in the risky asset. Consume nothing; maximise
expected utility of terminal wealth $U(X_T) = X_T^{1-\gamma}/(1-\gamma)$
(CRRA with coefficient $\gamma > 0$, $\gamma \ne 1$).

The wealth SDE is $dX_t = X_t[(r + \alpha_t(\mu-r))\, dt +
\alpha_t \sigma\, dW_t]$. The HJB equation is
\[
\partial_t V + \sup_\alpha\Big\{
x[r + \alpha(\mu-r)]\,\partial_x V
+ \tfrac{1}{2}\alpha^2 \sigma^2 x^2 \partial_{xx}V
\Big\} = 0.
\]
Optimising over $\alpha$ (take $\partial/\partial \alpha$):
\[
\alpha^*(x,t) = -\,\frac{(\mu-r)\,\partial_x V}{\sigma^2 x\, \partial_{xx}V}.
\]
For CRRA utility the value function separates as
$V(x,t) = f(t)\,x^{1-\gamma}/(1-\gamma)$, plugging in gives $\alpha^*$
constant:
\[
\alpha^* = \frac{\mu - r}{\gamma\,\sigma^2}.
\]
This is the \emph{Merton ratio}, identical to the Kelly fraction of
\cref{ch:kelly_risk} up to the risk-aversion scaling $\gamma$. The
log-utility limit $\gamma \to 1$ recovers full Kelly. The whole
apparatus of Avellaneda--Stoikov market-making (\cref{ch:market_making})
is an HJB in a richer state space (inventory plus cash), solved by the
same ansatz-and-reduce technique.

% =====================================================================
\section{Numerical tools}
\label{sec:app_math_numerics}

A handful of low-level numerical primitives appear throughout the
derivatives-pricing and calibration parts of the book. Each is sketched
here; full treatments with stability analysis and Python references are
in \cref{ch:numerical_methods}.

\emph{Bisection} finds a root of a continuous function on $[a, b]$
given $f(a)f(b) < 0$. Repeatedly halve the interval, keeping whichever
half still brackets a sign change. Linear convergence; one guaranteed
bit of accuracy per iteration; robust. Used for implied-volatility
inversion in \cref{ch:vol_surface} when a root bracket is easy.

\emph{Newton--Raphson} improves by using the derivative:
$x_{n+1} = x_n - f(x_n) / f'(x_n)$. Quadratic convergence near a simple
root. Fast but fragile: diverges for poor initial guesses or near
inflection points. The standard implied-vol solver uses Newton with a
safety bisection fallback (\cref{ch:numerical_methods}).

\emph{Worked Newton.} Solve $f(x) = x^2 - 2 = 0$ (so the true root is
$\sqrt{2} \approx 1.41421356$). With $f'(x) = 2x$ the iteration is
$x_{n+1} = \tfrac{1}{2}(x_n + 2/x_n)$. Starting from $x_0 = 1.5$:
\[
x_1 = 1.41666\ldots,\; x_2 = 1.41421568\ldots,\; x_3 = 1.41421356\ldots
\]
Three iterations, full double-precision accuracy. Error roughly squares
each step --- quadratic convergence in action. An implied-vol
Newton solve on a well-conditioned option starts from a Brenner--Subrahmanyam
initial guess and typically converges in 3--5 iterations.

\emph{Secant method}, a derivative-free cousin of Newton, estimates
$f'(x_n)$ from $(f(x_n) - f(x_{n-1}))/(x_n - x_{n-1})$. Superlinear
convergence (golden-ratio rate $\approx 1.618$), cheaper per iteration
when $f'$ is expensive. The default choice in SciPy's
\verb|brentq| / \verb|brenth| solvers combines it with bisection for
robustness.

\emph{Tridiagonal solve} (\emph{Thomas algorithm}) solves a system
$A\bm{x} = \bm{d}$ where $A$ is tridiagonal in $\mathcal{O}(n)$ time,
not $\mathcal{O}(n^3)$. This is what makes implicit finite-difference
schemes for parabolic PDEs (Crank--Nicolson for Black--Scholes) viable
at realistic grid sizes.

\emph{Monte Carlo basics}: estimate $\E[f(X)]$ by drawing $n$ i.i.d.\
samples $X^{(i)}$ and computing $\hat \mu_n = \frac{1}{n}\sum_i f(X^{(i)})$.
By the LLN this converges; by the CLT the standard error is
$\hat\sigma / \sqrt{n}$, giving the infamous
$\mathcal{O}(1/\sqrt{n})$ rate --- to halve the confidence interval
requires quadrupling the samples. Variance-reduction techniques
(antithetics, control variates, importance sampling, stratification)
are the only practical way to cut compute; see
\cref{ch:numerical_methods}.

\emph{Finite differences} approximate derivatives of a smooth function
by
\[
f'(x) \approx \frac{f(x+h) - f(x-h)}{2h} \quad (\text{central, } \mathcal{O}(h^2)),
\qquad
f''(x) \approx \frac{f(x+h) - 2f(x) + f(x-h)}{h^2}.
\]
Greeks (\cref{ch:greeks_hedging}) are computed this way when a
closed form is unavailable or cumbersome. Choice of bump size $h$
trades off truncation error (grows with $h$) against cancellation
error (grows as $h$ shrinks, as subtraction of nearly-equal numbers
loses significant digits); a double-precision sweet spot is typically
$h \sim 10^{-4}$ for first derivatives and $10^{-2}$ for seconds.

\emph{Optimisation.} Unconstrained minimisation of a differentiable
$f : \R^n \to \R$ reduces to solving $\nabla f = 0$. Standard methods:
gradient descent (simple, slow); Newton's method (quadratic convergence,
but requires Hessian); quasi-Newton (BFGS, L-BFGS: Hessian approximated
from gradient history, the workhorse of \verb|scipy.optimize.minimize|).
For convex problems every local minimum is global; most calibration
problems in quant finance (implied-vol fitting, SABR parameter fitting)
are non-convex and benefit from good initial guesses.

\begin{pitfallbox}[Do not invert matrices]
``Compute $A^{-1}\bm{b}$'' in code should almost never be written
\verb|np.linalg.inv(A) @ b|. Use \verb|np.linalg.solve(A, b)| (LU with
partial pivoting), or \verb|scipy.linalg.cho_solve| for PSD $A$, or
Thomas for tridiagonal $A$. Explicit inverses are both slower and
numerically worse.
\end{pitfallbox}

% =====================================================================
\section{Further reading}
\label{sec:app_math_refs}

The following are the standard references. Anything in this appendix
that was unfamiliar is best re-learned from one of them, not from the
internet.

\begin{itemize}
  \item \textbf{Probability.} \citet{grimmett2001probability},
        \emph{Probability and Random Processes} (Oxford). The standard
        undergraduate--graduate reference; clean proofs, wide coverage,
        good exercises. Chapters 1--7 cover everything in
        \cref{sec:app_math_prob} and more.
  \item \textbf{Linear algebra.} \citet{strang2016linalg},
        \emph{Introduction to Linear Algebra}. Concept-first, with the
        computational story never far from the theory. The classic
        complement is \citet{trefethen1997numerical},
        \emph{Numerical Linear Algebra}, for SVD, conditioning, and
        stable algorithms --- what to use in production.
  \item \textbf{Stochastic calculus.} \citet{oksendal2003stochastic},
        \emph{Stochastic Differential Equations}. The most readable
        graduate introduction; It\^o, SDEs, Feynman--Kac, and an
        extremely clean optional-stopping chapter.
        \citet{shreve2004stochastic} \emph{Stochastic Calculus for
        Finance II} is the finance-specialised alternative and is the
        text that sits on most derivatives strats' desks.
  \item \textbf{Stochastic control.} \citet{pham2009continuous},
        \emph{Continuous-time Stochastic Control and Optimisation with
        Financial Applications}. HJB, viscosity solutions, Merton-type
        problems. \citet{cartea2015algorithmic} applies the machinery
        directly to optimal market making and execution and is the
        single best companion to \cref{ch:market_making} and
        \cref{ch:almgren_chriss}.
  \item \textbf{Numerical methods for finance.}
        \citet{glasserman2004monte}, \emph{Monte Carlo Methods in
        Financial Engineering} --- the definitive reference for
        variance reduction and quasi-Monte Carlo in derivatives
        pricing. For PDE methods, \citet{duffy2006finite} is a
        standard desk reference.
\end{itemize}

\begin{keyfactbox}[Minimal prerequisite stack]
If starting from scratch and pressed for time, a viable one-semester
self-study path is: Strang (linalg, 8 weeks) $\to$ Grimmett--Stirzaker
chapters 1--5 (probability, 4 weeks) $\to$ \O ksendal chapters 1--5
(Brownian motion, It\^o, SDE, Feynman--Kac, 4 weeks). Everything else
in this appendix is either consolidation or an application of those
three books.
\end{keyfactbox}
